\chaptertopic{Introduction}

\begin{learningobjectives}
  \item Explain what a business process is.
  \item Identify common process characteristics.
  \item Explain why organizations model processes.
\end{learningobjectives}

\section{Business Processes}

A \term{business process} is a collection of related activities that produces
a specific outcome. Processes transform inputs into outputs and create value
for a customer or another stakeholder.

\definition[def:business-process]{Business Process}{
A structured collection of activities designed to achieve a particular
organizational objective. A process has an identifiable beginning, result,
and recipient of that result.

For a simple process, the transformation can be summarized as
\[
  \text{Input} + \text{Activities} \longrightarrow \text{Output}.
\]
}

\keyconcept{Process Owner}{
A process owner is accountable for the performance and ongoing improvement of a
particular business process. The role crosses organizational boundaries when
necessary to ensure the process meets its objectives.
}

\example{Online Order}{
An online purchase may consist of receiving an order, verifying payment,
preparing the package, shipping it, and confirming delivery. These activities
work together to produce the outcome: the customer receives the correct order.
}

\important{
A process should have a clearly identifiable input and output. Listing tasks
without identifying the result does not fully describe a process.
}

\note{
The people who perform a process do not necessarily own or manage it. Process
ownership is an accountability role, not simply a job title.
}

\section{A Simple Process View}

The following table gives an example of inputs, activities, and outputs.
\begin{center}
  \begin{tabularx}{\textwidth}{@{}lYY@{}}
    \toprule
    Stage & Example & Purpose \\
    \midrule
    Input & Customer order & Starts the process \\
    Activities & Validate, pack, ship & Transform the input \\
    Output & Delivered order & Meets the customer's need \\
    \bottomrule
  \end{tabularx}
\end{center}

\section{Revision Questions}

\begin{questionlist}
  \q[q:business-process]{What is a business process?}

  \q{What are the main characteristics of a business process?}

  \qa[q:process-owner]{What is a process owner?}
  {The person accountable for the performance and improvement of a specific
  business process.}

  \examquestion{Explain why defining a process output is important when
  deciding whether a collection of activities forms a process.}
\end{questionlist}

The definition is given in Definition~\ref{def:business-process}; Question~
\ref{q:business-process} and Question~\ref{q:process-owner} are useful for
revision.

\begin{summarybox}
The main ideas from this chapter are:
\begin{itemize}
  \item Business processes consist of related activities.
  \item Processes transform inputs into outputs.
  \item Every process should have a clear objective and outcome.
\end{itemize}
\end{summarybox}
